\section*{Chapter \ref{ch:b1:ecosystem-organization} --- Ecosystems and Trophic Structure}

\begin{solution}{exo:b1:ecosystem-organization:1}
The pond \omterm{def:b1:ecosystem-organization:ecosystem}{ecosystem} is the \omterm{def:b1:ecosystem-organization:ecosystem}{community} --- all its \omterm{def:b1:populations:population}{populations}: algae,
pondweed, snails, insect larvae, frogs, fish, bacteria --- together
with the \omterm{def:b1:ecosystem-organization:ecosystem}{biotope}: the water, its temperature, light, dissolved gases
and minerals, the mud.
\end{solution}

\begin{solution}{exo:b1:ecosystem-organization:2}
\omterm{def:b1:ecosystem-organization:niche}{Habitat}: the place a species lives (a spruce forest). \omterm{def:b1:ecosystem-organization:niche}{Niche}: its way
of living there --- what it eats, when, where in the tree, what it
tolerates. Fundamental \omterm{def:b1:ecosystem-organization:niche}{niche}: the range it could occupy alone;
realised: what competitors and predators leave it.
\end{solution}

\begin{solution}{exo:b1:ecosystem-organization:3}
$20\,800/1.7\times 10^6 = \qty{1.2}{\%}$; herbivores took
$3400/8800 = \qty{39}{\%}$ of net production.
\end{solution}

\begin{solution}{exo:b1:ecosystem-organization:4}
\omterm{def:b1:ecosystem-organization:trophic}{Producers}, \omterm{def:b1:ecosystem-organization:trophic}{consumers} (primary, secondary, tertiary), \omterm{def:b1:ecosystem-organization:trophic}{decomposers},
detritivores; the \omterm{def:b1:ecosystem-organization:trophic}{decomposers}, which return minerals from dead matter
to the \omterm{def:b1:ecosystem-organization:trophic}{producers}.
\end{solution}

\begin{solution}{exo:b1:ecosystem-organization:5}
NPP $= 800\times 17 = \qty{13.6}{MJ/m^2}$; eaten $\qty{2.04}{MJ}$;
assimilated $\qty{1.02}{MJ}$; rabbit $\qty{51}{kJ/m^2}$, i.e.\
\qty{510}{MJ} per hectare (\qty{30}{kg} of rabbit at \qty{17}{kJ/g}).
Efficiency $51/13\,600 = \qty{0.4}{\%}$.
\end{solution}

\begin{solution}{exo:b1:ecosystem-organization:6}
Biomass is a stock: phytoplankton that is eaten as fast as it grows
need never accumulate, so the standing crop of \omterm{def:b1:ecosystem-organization:trophic}{producers} can be less
than that of the longer-lived animals feeding on it. Energy is a
flow: every level receives all it will ever have from the one below
and must lose part of it as heat, so the flow can only shrink upward.
\end{solution}

\begin{solution}{exo:b1:ecosystem-organization:7}
$p = 0.5, 0.3, 0.15, 0.05$: $H = -(0.5\ln 0.5 + 0.3\ln 0.3 + 0.15\ln
0.15 + 0.05\ln 0.05) = 0.347 + 0.361 + 0.285 + 0.150 = 1.14$;
evenness $1.14/\ln 4 = 0.82$. Four at 25 each: $H = \ln 4 = 1.39$,
evenness 1.
\end{solution}

\begin{solution}{exo:b1:ecosystem-organization:8}
Net $= \qty{3}{mg}$ of $\mathrm{O_2}$ per litre per day; respiration
\qty{1}{mg}; gross $3 + 1 = \qty{4}{mg}$. In carbon: gross
\qty{1.5}{mg}, net \qty{1.1}{mg} of C per litre per day.
\end{solution}

\begin{solution}{exo:b1:ecosystem-organization:9}
Each link keeps a tenth: ten links would leave a billionth of the
primary production, too little to feed a single animal over any area.
Where production is low (desert) or the area small (island), the
tenth-of-a-tenth runs out after fewer links; a forest's high
production supports one or two more.
\end{solution}

\begin{solution}{exo:b1:ecosystem-organization:10}
Grain: \qty{1}{t} needs a sixth of a hectare. Beef: \qty{1}{t} of beef
energy needs \qty{10}{t} of grain, \num{1.7} hectares --- ten times
the land. Eating at the second level costs a level's worth of energy;
a planet feeds ten times more people on grain than on grain-fed meat,
and the human \omterm{def:b1:ecosystem-organization:trophic}{trophic level}, averaged over diets, is a large term in
the \omterm{thm:b1:populations:logistic}{carrying capacity}.
\end{solution}

\begin{solution}{exo:b1:ecosystem-organization:11}
Per square metre the ocean is a desert, but it covers two thirds of
the planet, so a low rate times an enormous area gives half the
total. Its production is limited by nutrients, which sink out of the
lit surface; where currents bring them back up (coastal upwellings,
shelves), production is ten times higher, and there, on a few percent
of the sea, the fish and the fisheries concentrate.
\end{solution}

\begin{solution}{exo:b1:ecosystem-organization:12}
Energy enters as light, passes through a few levels of \omterm{def:b1:organism-environment:organism}{organisms},
and \omterm{def:b1:flowering-plant-organization:organs}{leaves} entirely as heat: in that sense the \omterm{def:b1:ecosystem-organization:ecosystem}{ecosystem} does turn
sunlight into heat, and the \omterm{def:b1:ecosystem-organization:trophic}{decomposers}, respiring everything the
\omterm{def:b1:ecosystem-organization:trophic}{consumers} did not, are where most of it goes. But the matter is not
consumed: carbon, nitrogen and phosphorus cycle, returned by the same
\omterm{def:b1:ecosystem-organization:trophic}{decomposers} to the \omterm{def:b1:ecosystem-organization:trophic}{producers}, so the system persists; and the
\omterm{def:b1:organism-environment:organism}{organisms} are not a by-product but the structure through which the
flow is organised --- the \omterm{def:b1:ecosystem-organization:niche}{niches}, webs and pyramids are what the
energy builds on its way through. The sentence has the thermodynamics
and misses the biology.
\end{solution}

\begin{solution}{pb:b1:ecosystem-organization:1}
\textbf{1.} $20\,800 - 12\,000 = \qty{8800}{kcal}$ per square metre per
year.
\textbf{2.} \qty{1.2}{\%}. Losses: light not absorbed (half the
spectrum, reflection, water), \omterm{def:b1:flowering-plant-organization:organs}{leaves} saturated at full sun or shaded,
photorespiration, and the plants' own respiration (already counted
between gross and net).
\textbf{3.} $12\,000/20\,800 = \qty{58}{\%}$.
\textbf{4.} Eaten \qty{3400}{}; directly to \omterm{def:b1:ecosystem-organization:trophic}{decomposers}
$8800 - 3400 = \qty{5400}{kcal}$.
\textbf{5.} $8800/10 = \qty{880}{g/m^2}$: like a temperate forest or a
rich grassland.
\textbf{6.} $1.7\times 10^6 - 20\,800 = \qty{1.68e6}{kcal}$, \qty{98.8}{\%}
of it: reflected, transmitted, or absorbed by water and rock and
re-radiated as heat.
\textbf{7.} Production $1500 - 1100 = \qty{400}{kcal}$; assimilation
efficiency $1500/3400 = \qty{44}{\%}$.
\textbf{8.} $400/1500 = \qty{27}{\%}$.
\textbf{9.} $400/8800 = \qty{4.5}{\%}$ (\qty{16}{\%} if reckoned on what
was eaten, \qty{3400}{}).
\textbf{10.} Carnivore production $380 - 300 = \qty{80}{kcal}$;
efficiency $80/400 = \qty{20}{\%}$ (they ate \qty{380}{} of the
\qty{400}{} produced).
\textbf{11.} $6/80 = \qty{7.5}{\%}$.
\textbf{12.} \qty{6}{kcal}; $6/1.7\times 10^6 = \num{3.5e-6}$, a few
millionths.
\textbf{13.} Meat is digestible and close in composition to the eater;
plant matter is fibrous and largely indigestible, so a herbivore
loses more of its intake as faeces.
\textbf{14.} Uneaten net production \qty{5400}{}; herbivore faeces
$3400 - 1500 = \qty{1900}{}$; herbivore production not eaten
$400 - 380 = \qty{20}{}$; carnivore faeces (assimilation not given;
take eaten minus respired minus production $= 0$, all assimilated)
and uneaten carnivore production $80 - 20 = \qty{60}{}$; top-carnivore
production \qty{6}{}: total \qty{7386}{kcal}.
\textbf{15.} $12\,000 + 1100 + 300 + 14 + 7386 = \qty{20800}{kcal}$.
\textbf{16.} Equal to the gross production: the spring is in a steady
state, storing no biomass from year to year.
\textbf{17.} $7386/20\,800 = \qty{36}{\%}$ (most of the rest is the
\omterm{def:b1:ecosystem-organization:trophic}{producers}' own respiration).
\textbf{18.} \omterm{def:b1:ecosystem-organization:trophic}{Producers} \qty{8800}{}, herbivores \qty{400}{}, carnivores
\qty{80}{}, top carnivores \qty{6}{}; efficiencies \qty{4.5}{\%},
\qty{20}{\%}, \qty{7.5}{\%}.
\textbf{19.} \qty{8000}{}, \qty{400}{}, \qty{100}{}, \qty{15}{kcal/m^2}:
upright, each level a tenth or less of the one below.
\textbf{20.} \omterm{def:b1:ecosystem-organization:trophic}{Producers} $8000/8800 = 0.9$ years; herbivores $400/400 =
1$ year; carnivores $100/80 = 1.25$ years; top carnivores $15/6 =
2.5$ years. The plants turn over fastest: short-lived \omterm{def:b1:body-plans-tissues:tissue}{tissue}, grazed
continuously; the top carnivores are long-lived animals whose stock
is renewed slowly.
\textbf{21.} A heron of \qty{2}{kg} is about \qty{500}{g} dry,
\qty{5000}{kcal}, and needs about that in production each year to
replace itself and its young: $5000/6 \approx \qty{800}{m^2}$ of
spring --- in practice far more, since it also respires: at
\qty{100}{kcal} a day, \qty{36500}{kcal} a year of intake, i.e.\
\qty{6000}{m^2} of the level below it (carnivore production
\qty{80}{}), which is what a heron's territory looks like.
\textbf{22.} The extra net production is not eaten and goes to the
\omterm{def:b1:ecosystem-organization:trophic}{decomposers}, whose respiration doubles; at night, with no
photosynthesis, their oxygen demand can empty the water of oxygen
and kill the fish --- eutrophication.
\textbf{23.} Herbivore production falls to \qty{200}{}: carnivores get
half their food, their production falls toward \qty{40}{}, the top
carnivores' toward 3, and there are fewer herons; the eelgrass,
grazed less, thickens and more of it goes uneaten to the \omterm{def:b1:ecosystem-organization:trophic}{decomposers}.
\textbf{24.} The top carnivores' energy is tiny, but their effect is
on the carnivores' numbers, which control the herbivores, which
control the grass: removing them lets the carnivores rise, the
herbivores fall and the grass grow --- a cascade through the web
(\cref{ch:b1:species-interactions}); removing a few kilocalories of
herbivores changes nothing but a few kilocalories.
\textbf{25.} \omterm{def:b1:ecosystem-organization:trophic}{Producers} to herbivores \qty{4.5}{\%} (of NPP; \qty{16}{\%}
of what was eaten), herbivores to carnivores \qty{20}{\%}, carnivores
to top carnivores \qty{7.5}{\%}; sunlight to gross production
\qty{1.2}{\%}; a few millionths of the sunlight ends in the top
carnivores.
\end{solution}
